% =========================================================================
% GUÍA DE CONFIGURACIÓN Y ESTILOS DE LA PLANTILLA SOBRIA
% Diseñada con el formato Charter para Ingeniería y Ciencias.
% =========================================================================
\documentclass[11pt, a4paper]{article}

% --- Paquetes de Estilo ---
\usepackage[utf8]{inputenc}
\usepackage[T1]{fontenc}
\usepackage[spanish, es-tabla]{babel}
\usepackage{amsmath, amssymb, amsfonts}
\usepackage{graphicx}
\usepackage{booktabs}
\usepackage{xcolor}
\usepackage{listings}
\usepackage{charter}
\usepackage{microtype}
\usepackage{metalogo}      % Corrige el espaciado y kerning de los logos (ej. \LaTeX) con fuentes personalizadas
\usepackage{setspace}
\setstretch{1.12}

% --- Geometría ---
\usepackage[margin=2.5cm]{geometry}

% --- Colores ---
\definecolor{primary}{rgb}{0.0, 0.35, 0.35}  % Verde azulado sobrio
\definecolor{secondary}{rgb}{0.2, 0.2, 0.2} % Gris carbón
\definecolor{codebg}{rgb}{0.97, 0.97, 0.95}
\definecolor{grayborder}{rgb}{0.85, 0.85, 0.85}

% Aplicar color secundario al cuerpo del texto
\makeatletter
\newcommand{\globalcolor}[1]{%
  \color{#1}\global\let\default@color\current@color
}
\makeatother
\AtBeginDocument{\globalcolor{secondary}}

% --- Hipervínculos ---
\usepackage{hyperref}
\hypersetup{
    colorlinks=true,
    linkcolor=primary,
    urlcolor=primary,
    citecolor=primary,
    pdftitle={Guía de Configuración: Plantilla Sobria},
    pdfauthor={Jair Aguilar Peña}
}

% --- Listings ---
\lstdefinestyle{customcode}{
    backgroundcolor=\color{codebg},
    commentstyle=\color{primary!70}\itshape,
    keywordstyle=\color{primary}\bfseries,
    numberstyle=\tiny\color{gray},
    stringstyle=\color{orange!80!black},
    basicstyle=\ttfamily\small,
    breakatwhitespace=false,
    breaklines=true,
    captionpos=b,
    keepspaces=true,
    numbers=left,
    numbersep=5pt,
    frame=single,
    rulecolor=\color{grayborder},
    aboveskip=15pt,
    belowskip=10pt
}
\lstset{style=customcode}

% =========================================================================
% DATOS DEL DOCUMENTO
% =========================================================================
\title{\textbf{\color{primary}Guía de Configuración: Plantilla Académica Sobria}}
\author{Jair Aguilar Peña}
\date{\today}

% =========================================================================
% DOCUMENTO
% =========================================================================
\begin{document}

\maketitle

\begin{abstract}
Este documento describe de manera detallada las decisiones de diseño, paquetes utilizados y configuraciones que componen la plantilla sobria y minimalista orientada a ciencias e ingeniería. El objetivo es explicar el porqué de cada elemento y cómo este formato asegura un aspecto profesional y legible, listo para cumplir con rigurosos filtros de publicación.
\end{abstract}

\section{Introducción}
El diseño de documentos científicos requiere un balance entre estética y apego a normas rigurosas. Tradicionalmente, los documentos de \LaTeX\ pecan de verse demasiado densos o visualmente anticuados por usar la tipografía predeterminada \textit{Computer Modern} sin configuraciones adicionales de espaciado.

Esta plantilla propone resolver ese problema usando tres principios:
\begin{enumerate}
    \item \textbf{Legibilidad de pantalla:} El $95\%$ de las revisiones académicas se hacen en pantallas de computadora o tabletas.
    \item \textbf{Espaciado descansado:} Suficiente aire en márgenes y entre líneas sin perder densidad de información.
    \item \textbf{Adaptabilidad al cambio:} Configuración puramente semántica que permite migrar a formatos como IEEEtran o ACM con mínimos cambios.
\end{enumerate}

\section{Análisis de la Configuración y Paquetes Clave}
A continuación se detalla el propósito de cada paquete seleccionado en la plantilla.

\subsection{Tipografía Charter (\texttt{\textbackslash usepackage\{charter\}})}
Diseñada por Matthew Carter en 1987 para impresoras de baja resolución de la época, la fuente \textit{Charter} posee caracteres con serifas robustas y formas abiertas. 
\begin{itemize}
    \item \textbf{Por qué se usa:} Sus trazos gruesos y equilibrados se renderizan extremadamente bien en pantallas modernas de alta resolución, evitando la delgadez que a veces hace que la tipografía tradicional se vea "borrosa" o difícil de leer en pantallas medianas.
    \item \textbf{Compatibilidad:} Es ideal para ingeniería, ya que combina perfectamente con símbolos matemáticos y denota seriedad sin la rigidez de \textit{Times New Roman}.
\end{itemize}

\subsection{Microtipografía (\texttt{\textbackslash usepackage\{microtype\}})}
Este paquete realiza ajustes dinámicos invisibles a nivel de caracteres (protrusión y expansión de fuentes).
\begin{itemize}
    \item \textbf{Por qué se usa:} Elimina las molestas "calles" o espacios en blanco gigantescos que ocurren cuando \LaTeX\ intenta justificar a la fuerza un párrafo. Hace que el bloque de texto se vea uniforme y pulcro, elevando drásticamente el nivel estético del documento.
\end{itemize}

\subsection{Interlineado Moderno (\texttt{\textbackslash setstretch\{1.12\}})}
Por defecto, \LaTeX\ compila con un interlineado muy cerrado ($1.0$), diseñado originalmente para ahorrar papel de imprenta.
\begin{itemize}
    \item \textbf{Por qué se usa:} Un interlineado de $1.12$ a $1.15$ (habilitado mediante el paquete \texttt{setspace}) introduce el aire suficiente entre líneas para descansar la vista en lecturas largas, sin cruzar el límite hacia el "doble espacio" que suele verse desordenado en documentos cortos.
\end{itemize}

\subsection{Geometría de Página (\texttt{\textbackslash usepackage[margin=2.5cm]\{geometry\}})}
Los márgenes predeterminados de \LaTeX\ son extremadamente anchos.
\begin{itemize}
    \item \textbf{Por qué se usa:} Establecer los márgenes en $2.5\text{ cm}$ optimiza el uso del espacio en papel A4 o Carta. Es el tamaño estándar aceptado por la mayoría de las universidades y reportes oficiales.
\end{itemize}

\subsection{Paleta de Colores (\texttt{xcolor})}
El minimalismo radica en usar color con un propósito funcional, no decorativo.
\begin{itemize}
    \item \textbf{\texttt{primary} (Teal obscuro):} Se utiliza para títulos y enlaces web (`href`). Aporta un toque distintivo sin la agresividad del rojo o el azul eléctrico predeterminado de \LaTeX.
    \item \textbf{\texttt{secondary} (Gris carbón):} En lugar de negro puro (`\#000`), el cuerpo del texto está tintado en un gris muy obscuro (`rgb(0.2, 0.2, 0.2)`). Esto suaviza el contraste en pantallas retroiluminadas y previene la fatiga ocular.
\end{itemize}

\section{Buenas Prácticas de Escritura Semántica}
Para asegurar que tu documento pase los filtros de verificación oficiales sin contratiempos, sigue estas reglas de estructuración:

\subsection{Tablas con \texttt{booktabs}}
El uso de líneas verticales (ej. \texttt{|c|c|}) y rejillas completas es considerado mala práctica en publicaciones científicas modernas. En su lugar, usa el paquete \texttt{booktabs} y limita tu diseño a tres líneas horizontales clave:
\begin{lstlisting}[language=TeX, caption={Sintaxis correcta de tablas académicas.}]
\begin{table}[h]
\centering
\caption{Título de la tabla (SIEMPRE ARRIBA).}
\begin{tabular}{lll}
    \toprule
    Cabecera 1 & Cabecera 2 & Cabecera 3 \\
    \midrule
    Dato A     & Dato B     & Dato C \\
    Dato D     & Dato E     & Dato F \\
    \bottomrule
\end{tabular}
\end{table}
\end{lstlisting}

\subsection{Imágenes y Código}
Recuerda que para imágenes (`figure`) y listados de código (`lstlisting`), la descripción (o caption) debe ir **siempre debajo** del elemento. El lector necesita procesar la información visual primero antes de leer su descripción detallada.

\end{document}
